\pagebreak

{
\thispagestyle{empty}

\mbox{}

\vspace{5cm}

\footnotesize
\noindent{}\emph{\noindent{}Ao final da tarde do sexto dia da criação divina, \emph{erev shabat}, às vésperas do dia sagrado, de descanso religioso e espiritual, alguns minutos antes do pôr do sol, Rabi Hanina e Rabi Oshea, dois rabinos talmúdicos do século \textsc{iii}, se reuniam nas encostas do Monte Moriá,\footnote{O Monte Moriá é um território sagrado judaico e o local escolhido para a construção do primeiro e do segundo Templo Sagrado de Jerusalém. Possui um simbolismo cósmico e testemunhou uma das máximas práticas de benevolência judaica, \emph{veahavtá lereachá kamocha}, ou seja, ``ame o teu próximo como a ti mesmo''. Sobre isso, diz Rabi Akiva: \emph{Ze'klal gadol baTorá}, ou seja, ``este é o maior princípio da Torá''. A filosofia judaica acredita que o terceiro Templo Sagrado de Jerusalém, bem como a salvação divina, virá até o ano 6\,000 do calendário judaico. Atualmente nos encontramos no ano 5\,787.} em Jerusalém. Nesse recinto sagrado, moldavam e esculpiam conjuntamente um animal em barro e gravavam em sua testa a palavra \emph{emet}.\footnote{\emph{Emet}, além de significar em português ``verdade'', é um dos nomes judaicos de Deus.} Com o auxílio do \emph{Sêfer Ietzirá},\footnote{O \emph{Sêfer Ietzirá}, o ``Livro da Criação'', é um dos mais antigos tratados filosóficos escritos em língua hebraica.} traziam à vida um bezerro e, logo em seguida, o cozinhavam como alimento para a ceia festiva. Em \emph{Terra e paz}, o escritor e poeta israelense Yehuda Amichai nos lembra que a procura por um bezerro\footnote{No original, um cabrito. Ou qualquer outro animal de pasto.}, ou por um filho\footnote{O Monte Moriá foi o local em que o primeiro patriarca Abraão foi impelido a sacrificar seu filho Isaac a fim de provar a sua lealdade a Deus.} no Monte Moriá sempre foi o começo de uma nova religião.\footnote{\textsc{amichai}, \emph{Terra e Paz}, p.\,29.}}

\pagebreak
}

\chapter{Páginas}

% Comentário da editora: Gostei desse prelúdio. Podemos pensar um para cada capítulo.

% \begin{center}\adforn{49}\end{center}

\noindent{}Conta o Talmud que, durante a época do segundo Templo Sagrado de
Jerusalém, uma jovem mulher, proibida de transitar por esta
territorialidade\footnote{Na maioria das comunidades judaicas, e em
  especial naquelas comunidades tradicionais, ortodoxas ou
  ultraortodoxas, mulheres ainda são proibidas de assumir papéis de
  liderança no âmbito religioso e comunitário.}, envolve-se em uma
conversa com Deus, por meio de palavras. A jovem suplica a Deus para que
conceda a ela um filho para criar, permitindo, desta maneira, que seja
portadora de uma nova linhagem de descendentes e matriarca destes fios
genealógicos. Assim, ao envolver-se em uma conversa com Deus, a jovem
mulher passa a ser, ela mesma, porta-voz de sua identidade, de sua
conexão e de sua relação pessoal com a divindade\footnote{No judaísmo
  pós-templo, qualquer pessoa pode se conectar com Deus e com a
  divindade sem a necessidade de uma intermediação e encontrar
  interpretações próprias nos textos sagrados.}.

Entre 586 a.E.C. e 76 d.E.C., período em que perdurou o segundo Templo
Sagrado de Jerusalém, os hebreus se comunicavam com a divindade por meio
de sacrifícios e oferendas realizadas e concedidas pelo \emph{Cohen
Hagadol} --- o Sumo Sacerdote de Israel. Este era o canal da comunicação
direta e unilateral com Deus.

Após a destruição do Templo Sagrado, a expulsão da
\emph{Shechiná}\footnote{\emph{Shechiná} é uma palavra hebraica que, de
  maneira resumida, significa ``habitação'' ou ``presença de Deus''.}, e
assim a inexistência e a insuficiência de um território próprio judaico
para a realização de oferendas e de sacrifícios, os judeus passam a
comunicar-se com a divindade por meio de livros, de textos e do
pergaminho. Deus e, em especial, a memória do Templo Sagrado\footnote{A
  memória do Templo Sagrado e a lembrança da cidade de Jerusalém
  permeiam as rezas, os ritos, as celebrações, as festividades, bem como
  os momentos de luto judaico. Todas as sinagogas devem ser construídas
  mirando o Templo Sagrado (uma Jerusalém espiritual), voltando as
  orações a esta localidade santa. Nos noivados judaicos há o costume
  das mães dos noivos quebrarem conjuntamente um prato, nos casamentos
  judaicos quebra-se um copo, nas casas judaicas há uma parede sem
  revestimento (apenas em tijolos). Todas essas ações remetem à
  incompletude e à ausência do Templo Sagrado.} passam, dessa forma, a
habitar o texto e o livro.

Para manter a continuidade judaica viva, este novo aspecto da
comunicação com a divindade é santificado por meio de rezas e de cultos
religiosos entoados coletivamente, nas primeiras sinagogas da diáspora.
As rezas e as orações judaicas são conversas com Deus. São poemas,
cânticos e bênçãos que reiteram a lealdade e a fidelidade ao monoteísmo
e à divindade judaica. A reza judaica é um sopro entoado e proferido em
voz alta (em hebraico) e em comunidades em que os indivíduos judeus se
unem --- formando um \emph{minian}\footnote{\emph{Minian} é o quórum
  mínimo de dez homens judeus adultos (pós-\emph{bar-mitzvá}, maioridade
  religiosa) necessários para realizar certas obrigações religiosas, bem
  como para a retirada da Torá da Arca Sagrada e sua leitura.
  Mais recentemente e em algumas vertentes religiosas e liberais
  judaicas, as mulheres também passaram a contar como participantes
  ativas do \emph{minian}.} --- para realizar os ritos, para rezar e
santificar o nome de Deus de maneira coletiva.

Segundo a narrativa bíblica da criação Deus moldou, à sua imagem, o
homem em barro (em hebraico, \emph{Adamá}), o qual insuflou com o Seu
sopro. \emph{Adam} (Adão), o primeiro ser humano, nasce do cruzamento da
terra, da mão de Deus que confeccionou e moldou o barro e do sopro
divino --- adicionando vida, espírito e uma alma a esta matéria amorfa.

Conta o Talmud que Deus criou o mundo a partir do ar primal e
gravou, nesta esfera invisível, as 22 letras disponíveis do alfabeto
hebraico. É por meio deste gesto que todas as coisas do mundo, todos os
seres dos três estratos do cosmo --- espaço, tempo e o corpo humano ---
passam a existir. Assim, ``as palavras existem porque correspondem
a coisas existentes, e as coisas deveriam existir porque há palavras
para nomeá-las''\footnote{\textsc{manguel}, \emph{A musa da impossibilidade}, p.\,42.}. Por fim, o encadeamento linear das 22 letras hebraicas se torna
palavras que nomeiam coisas e, finalmente, compõem frases e orações.

No momento da Inquisição, entre as entranhas da espadana do resquício do
fogo, em que as inscrições das palavras sagradas se desprendem, voando
em direção aos céus, uma imagem assemelha-se ao ar primal de Deus. Tais
letras, ao se desgrudarem dos papéis, superam a própria morte, pousam em
outros territórios, criam mundos e fixam-se nestes espaços de vácuo. Ou
seja, as letras, ao morreram para renascer, tornam estes espaços lugares
de apropriações comunitárias e reiteram a criação do mundo por Deus.

Assim como o verbo ``escrever'' (em latim, \emph{scribere}), um gesto de
fazer uma incisão sobre um objeto e/\,ou corpo, o homem foi inscrito e
criado a partir do barro, ou seja, gravado e inserido nele. Ao final,
``os seres humanos'', como menciona Flusser, ``são
materiais espiritualizados da fusão da terra com Deus''\footnote{\textsc{flusser},
  \emph{A escrita. Há futuro para a escrita?}, p.\,25.}. Ao inscrever o
homem no barro, Deus grava, fura e penetra uma materialidade --- a molda
---, se apropria de sua maleabilidade física para criar o homem e,
assim, trazê-lo à vida.

Deus insere nesta matéria amorfa e desprovida de inteligência própria
--- um \emph{golem} --- uma alma por meio do sopro. Talvez essa seja
também a ação realizada pela jovem mulher. Ao rezar, ao iniciar um
diálogo com Deus, a jovem sopra palavras que são proferidas, voam e
vagueiam pelo ar até ganharem vida --- uma alma ---, até serem
espiritualizadas. Essas palavras, na forma de orações, flutuam em
direção aos céus, rompem os portais celestiais divinos, permitindo que a
divindade seja finalmente acessada sem uma intermediação.

Como diz a escritora e poeta brasileira Marília Garcia, ``as falas
acontecem no presente''. As falas e orações que saem de nossas bocas
por meio de vultos e frestas entre a língua e o ar da linguagem são
todas entoadas e ressoadas no momento do agora.

As falas, as palavras, as comunicações, bem como as histórias
testemunhadas e transmitidas oralmente (ou de forma escrita) de geração
em geração, acontecem no momento presente e, assim, se adiciona,
acrescenta e multiplica a alma, dando vida às comunidades, aos ritos e
aos cultos judaicos. Estes, por fim, também são interpretados de modos
diferentes, segundo o momento histórico, adicionando camadas de
complexidade às tradições.

Ao soprar, Deus inseriu uma alma e um espírito nos primeiros seres
humanos --- Adão e Eva --- e os contemplou com a vida. Ao final,
``o espírito penetra o objeto por ocasião do inscreve, para \emph{dar-lhe-vida}''\footnote{\emph{Ibidem}, p.\,25.}. Dessa forma, ao inscrever
uma informação (uma alma) em uma matéria amorfa, o autor --- Deus na
criação do mundo e de todas as coisas que ele contém, ou os seres
humanos por meio de frases e de palavras\footnote{Posso não acreditar em
  Deus, mas acredito no poder das palavras --- e, para mim, é quase a
  mesma coisa.} --- deseja que estes objetos e/\,ou corpos, por ele
inscritos, ressoem, perdurem no espaço e no tempo e, ao menos, não se
decomponham tão rapidamente.

% Comentário da editora: Tem algo errado na citação acima.

Estes objetos e/\,ou corpos inscritos e escritos --- possivelmente lapsos
de fagulhas ---, ou melhor, a vontade de perdurar e perfurar o espaço e
o tempo, se fazem sentir nos papéis do grupo judaico \emph{Oyneg Shabes}
como um nítido desejo de vida e de persistência. E assim também nos
escritos de minha trisavó Dvora. Uma vontade de adicionar --- a todo
custo --- a estas folhas, a estas páginas e a estes documentos (e a
estas palavras) uma vida própria, uma alma e um espírito que fure o
tempo, para que não nos esqueçamos tão rapidamente.

Em \emph{Teste de Resistores}\footnote{GARCIA, \emph{Teste de
  Resistores}, p.\,14.}, Marília Garcia comenta o diálogo entre Ferdinand
e Marianne em \emph{Pierrot le Fou}, filme de Jean-Luc Godard. Após um
assassinato e o roubo de um carro, Ferdinand quer parar para descansar
em uma praia. Ele se diz apaixonado, mas Marianne parece pensar apenas
em dinheiro.

\begin{verse}
nesse momento ferdinand se vira\\
olhando para trás na direção da câmera\\
e diz --- estão vendo\\
ela só pensa em se divertir\\
marianne se vira também olhando para trás\\
na direção da câmera\\
e pergunta para ele --- com quem você está falando?\\
ao que ele responde --- com o espectador.
\end{verse}

Este pequeno diálogo contribui para dar ao filme a dimensão de filme ---
apresentando o furo do cinema. Há pessoas que acreditam que não poderiam
viver sem escrever, \emph{Oyneg Shabes} e minha trisavó Dvora escrevem e
``quem escreve ultrapassa o ponto final em direção ao outro, o
leitor''\footnote{\textsc{costa} \emph{apud} \textsc{flusser}, \emph{op.\,cit.}, p.\,8.}. Essas pessoas
acreditam que o gesto de escrever pode expressar e atestar as suas
existências\footnote{Eu também escrevo porque as palavras me pertencem,
  e assim reafirmo a minha fé no poder da linguagem.}. Assim,
\emph{Oyneg Shabes} e minha trisavó estendem a mão em direção ao outro a
fim de provocar uma inquietação; comunicam-se com seus interlocutores e
permitem que suas palavras escritas falem por si.

Por fim, supõe-se que a inscrição (a criação) feita por Deus, dos
primeiros seres humanos, é repetida simbolicamente todos os anos, já que
todos os anos o mundo é novamente cosmogonizado por Deus e a aliança é
reafirmada. Segundo a filosofia judaica, Deus, no momento de \emph{Iom
Kipur}\footnote{O \emph{Iom Kipur} é uma das mais importantes
  festividades do calendário judaico, momento em que a proximidade de
  Deus permite a purificação dos pecados diante da divindade e com o
  próximo.}, inscreve e sela (outra vez) a humanidade em mais um ano de
vida, reitera a confiança na criação divina do mundo, e assim expande
novamente o sopro, a alma e uma vida.

Essa afirmação da aliança, e da lealdade judaica com Deus por meio da
palavra, é também realizada durante o momento do \emph{brit-milá} (em
português, ``aliança da circuncisão''). O \emph{brit-milá} é uma
cerimônia ritualística em que o prepúcio de meninos no oitavo dia de
vida é cortado, como símbolo da lealdade entre Deus e o povo judeu. É
neste momento, também, que o menino judeu recebe o seu nome diante da
comunidade --- e, assim, passa a existir perante ela e a frequentá-la.

Em hebraico, \emph{milá} --- a depender de sua forma escrita --- pode
significar tanto ``circuncisão'' ({מילה}), quanto ``palavra''
({מלה}); dessa maneira, uma tradução possível deste ritual (que
envolve um recém-nascido --- uma nova geração de indivíduos judeus)
seria: ``aliança da palavra''. Assim, a lealdade judaica a Deus e a
crença nas palavras são reafirmadas nos corpos judaicos por meio de uma
aliança com a lei judaica escrita, com a Torá, com os textos, com
as tradições e com a divindade --- e com a comunidade.

Se as falas, as rezas e as bênçãos são entoadas enquanto sopros, e
estes, por sua vez, acionam uma alma e um espírito, então, o velho vento
do sopro nômade que a diáspora carrega são as histórias e as palavras
esparsas, fagulhas e chispas de papéis. Territorialidades de trânsito
passíveis de apropriação e de pertencimentos postas em palavras. Textos
que são as chamas de vida e de celebrações judaicas religiosas e
diaspóricas. E que venham ventos mais leves.

Por fim, a \emph{Shechiná}, a presença e a habitação da divindade, que
deixou o Templo Sagrado (e que passou a habitar as palavras, os textos e
as interpretações), foi espalhada, peneirada, penetrada e permeada pelo
mundo e pela diáspora. Cabe a nós buscá-la, até que estes sentidos se
desvaneçam e sejam refeitos em outros territórios.

\begin{quote}
Os ventos ensinam, como diz um velho provérbio dos congos, que os
pássaros têm asas porque elas lhes foram passadas por outros
pássaros\footnote{\textsc{simas}; \textsc{rufino}, \emph{Encantamento. Sobre política de
  vida}, p.\,15.}.
\end{quote}

\begin{center}\adforn{49}\end{center}

\noindent{}Escrever é riscar as páginas. Este estilo riscante, como nos lembra
Vilém Flusser, se apresenta enquanto um estilete que dilacera e
esfarrapa imagens, ao final, ``as inscrições são imagens
esfarrapadas e rasgadas, imagens que o mortífero estilete
sacrificou''\footnote{\textsc{flusser}, \emph{op.\,cit.}, p.\,30.}. O estilete risca e
penetra a pedra ou o mármore para escrever, gravar nele suas breves
passagens, suas memórias, seus anseios e seus desejos de adeus e de
amor. No aramaico e no hebraico (ou no árabe), o movimento do estilete
sobre a pedra --- como nas \emph{Tábuas da Lei} --- é feito da esquerda
para a direita, em direção ao coração\footnote{Em uma tradução nunca se
  traduz palavras, mas sim sentimentos. Martin Buber nos conta que:
``O Rabi Mendel cuidava de que seus \emph{hassidim} não usassem nada no
  pescoço durante a oração. Pois, segundo dizia, não deve haver
  interrupção entre o cérebro e o coração''. \textsc{buber}, \emph{op.\,cit.}, p.\,608.}.

A escrita, o ato de gravar e/\,ou depositar tinta sobre uma superfície,
também diz respeito à visualidade. Em ``Metaescrita'', Vilém Flusser usa
um feliz exemplo: ``\emph{palavra} é uma palavra, mas \emph{frase} não é
uma \emph{frase}. Só é possível reconhecer isso na escrita, pois uma sentença
como essa, se pronunciada, soa incoerente''\footnote{\textsc{flusser}, \emph{op.\,cit}., p.\,20.}.

Assim, ``escrever é afirmar o paradoxo de que, como no gesto
divino, as palavras criam o mundo mas não dão conta de aprendê-lo no
livro''\footnote{\textsc{manguel}, \emph{op.\,cit.}, p.\,33.}. Em \emph{O Golem de
Praga}, livro originalmente em iídiche do Prêmio Nobel de
Literatura (1978) Isaac Bashevis Singer, urge uma tentativa de criação
de mundos por meio da escrita e da palavra.

Em meados do século \textsc{xviii}, o \emph{Maharal} (Grão-Rabino) de Praga, Rabi
Loeb, confecciona, a partir do barro, uma matéria amorfa e sem alma ---
um \emph{golem} (um ``boneco'') ---, com o intuito de ajudar a
comunidade judaica local perseguida pelos \emph{pogroms}\footnote{\emph{Pogrom}
  (``massacre'') é uma palavra russa que se refere às perseguições
  deliberadas e aos atos em massa de violência, espontânea ou
  premeditada, contra as comunidades judaicas no Leste Europeu.}. Sobre
a testa da criatura, o Rabi inscreve a palavra hebraica \emph{emet} (em
português, ``Verdade'') e assim a figura amorfa --- proveniente da terra
--- ganha vida, ajudando a comunidade nas tarefas diárias. Mas o boneco
logo escapa às instruções do Rabi, obrigando o \emph{Maharal} a
devolvê-lo ao pó. Para isso, Rabi Loeb apaga a primeira letra da palavra
hebraica \emph{emet} ({אמת}), transformando-a em \emph{met}
({מת}), em português: ``morto''. O boneco então desmancha-se ---
retornando ao barro\footnote{Rabi de Guer costumava dizer que o homem
  foi posto no mundo para que transformasse o pó em espírito.}.

Dessa forma, ``embora a primeira letra da inscrição em sua testa
tenha sido apagada e \emph{emet} tenha virado \emph{met}, uma palavra continua ao
nosso dispor para nomear mais um inominável: a própria
morte''\footnote{\textsc{manguel}, \emph{op.\,cit.}, p.\,45.}. E não seria a morte o
objetivo de toda a criação? E não seria Deus ``aquele que renova
diariamente a obra da criação''\footnote{\textsc{buber}, \emph{op.\,cit.}, p.\,647.}?

A partir da presente narrativa, entende-se que, acima de tudo, o
artifício mágico da palavra cria, descria e recria mundos, assim como os
destrói --- para então renascer. A palavra cria o \emph{golem} --- ou o
ser humano, por meio da palavra divina ---, atesta a sua existência, o
nomeia, como também reitera a sua insignificância ao inscrevê-lo
novamente, ao destruí-lo.

A bênção judaica \emph{shehakol nihiyah bed'varo} (em português:
``por cuja palavra todas as coisas vieram a existir'' ou ``que tudo
tenha sido feito a Sua palavra'') é entoada sobre todos e quaisquer
alimentos (antes de seu consumo) que não provenham diretamente da terra,
por exemplo: peixes, carnes, frangos, ovos, queijos, mel e todas as
bebidas (exceto vinho) --- além de comidas consideradas mágicas, tais
quais os cogumelos e o chocolate.

Se Deus moldou o mundo --- e todas as coisas que ele contém --- por meio
do cruzamento entre a terra e a Sua mão (que confeccionou o barro),
então, todas as demais coisas que não são diretamente provenientes da
terra --- como os animais, embora vivam nela --- foram, necessariamente,
feitas à Sua palavra --- por meio do sopro.

Há algo curioso ao analisarmos o início do texto bíblico. Se partirmos
do pressuposto de que \emph{Berechit Bará Elohim et Hashamaim ve et
HaHaaretz}\footnote{\emph{Berechit Bará Elohim et Hashamaim ve et
  HaHaaretz} é a primeira frase da narrativa bíblica, traduzida para o
  português como: ``No princípio, Deus criou o céu e a Terra''.} e o
conectivo \emph{et} ({את}) entre as palavras \emph{Elohim} e
\emph{Hashamaim} seria desnecessário e está escrito com a primeira
({א}) e a última letra ({ת}) do alfabeto hebraico, então Deus
criou, antes de mais nada, tudo que está entre a primeira e a última
letra, e por fim, \emph{Hashamaim} --- o céu --- \emph{ve et HaHaaretz}
--- e a terra. Deus finalmente nomeou todas as coisas que existem, e
assim todas as coisas passaram a existir por conta de suas palavras e de
seus nomes, já que, segundo Alberto Manguel, ``é possível
articular em sílabas tudo que existe''\footnote{\textsc{manguel}, \emph{op.\,cit.}, p.
  35.}.

% Comentário da editora: Normalmente é escrito com "s", "Bereshit". Os irmãos Campos, no livro da Perspectiva, traduzem como "Bere'shith". Enfim, o hebraico do livro inteiro precisa ser revisto e padronizado (não há uma norma acordada, mas eu tenho a minha própria, com base no que tem sido feito). Vale também para a acentuação conforme fonética portuguesa.

\begin{center}\adforn{49}\end{center}

% Comentário da editora: Essa história ficou meio jogada aqui. Vamos trabalhar com esse conceito ou não?

\noindent{}Certa vez, ``perguntaram ao Rabi Baruch: Por que dizemos: Bendito
o que falou, e o mundo se fez, e não Bendito o que criou o mundo? Ele
respondeu: Louvamos a Deus porque criou nosso mundo com palavras, e
não, como todos os outros mundos, com o pensamento''\footnote{\textsc{buber},
  \emph{op.\,cit.}, p.\,131.}.

\begin{center}\adforn{49}\end{center}

\noindent{}Na história bíblica de Babel, a humanidade tentou chegar a Deus por meio
da construção de uma torre que alcançasse os céus. Manguel, a partir da
leitura de \emph{Gênesis}, nos conta que, neste cenário, Deus
``fragmentou o idioma único que toda a humanidade falava até o
momento na miríade de idiomas que falamos hoje''\footnote{\textsc{manguel}, op.
  cit., p.\,35.} --- o que inclui todas as línguas, vivas ou mortas. Em
contrapartida o Rabi Pinkhas sustenta que, ``antes da construção
da Torre, todos os povos tinham em comum a língua sagrada, mas além
dela, cada um tinha também o seu próprio idioma. Por isso é que está
escrito: \emph{Toda a terra tinha uma língua}, que era a sagrada, \emph{e alguns
falares}, que são as línguas adicionais e separadas de cada povo. Nestas
cada povo se entendia entre si, naquela os povos se entendiam entre si.
O que Deus fez, ao castigá-los, foi tirar a língua
sagrada''\footnote{\textsc{buber}, \emph{op.\,cit.}, p.\,177.}.

No conto ``O Congresso'', do escritor Jorge Luis Borges, ``um
homem sonha em compilar uma enciclopédia completa do mundo e no fim se
dá conta de que a enciclopédia em si já existe --- e é o próprio
mundo''\footnote{\textsc{manguel}, \emph{op.\,cit.}, p.\,43.}. Assim, como ocorreu na
história de \emph{Babel}, é impossível alcançar Deus porque Deus, em si,
é o próprio mundo.

Para a \emph{Cabalá}, Deus possui 72 nomes apresentados em três
versículos codificados. HaMakom e HaShem são apenas dois deles. HaShem,
em português, significa ``o Nome''. Já HaMakom significa ``o Lugar''.
Assim sendo, Deus é o Nome e o Lugar em que tudo ocorre, no qual existe
tudo o que existe. Deus abriu um espaço em si para a criação do mundo
--- e de todas as coisas que ele contém --- bem como da humanidade, ou
seja, ``a realidade de Deus e a representação que Deus faz da
realidade são idênticas''\footnote{\emph{Ibidem}.}.

Martin Buber nos conta que, certa vez, ``perguntaram ao Rabi
Pinkhas: Por que é que Deus é chamado de Lugar? Ele é certamente o
Lugar no Universo, mas deveriam então chamá-lo assim, e não apenas
Lugar? Ele responde: O homem deve entrar em Deus, de modo que Deus
possa envolvê-lo e tornar-se o seu Lugar''\footnote{\textsc{buber}, \emph{op.\,cit.}, p.\,167.}.

Se Deus é o Lugar em que tudo acontece e se materializa, então, a
escrita e as palavras materializam o texto no lugar. Pois ``há um
momento no qual descobrimos que das marcas de tinta na página emerge um
mundo completamente formado e magicamente real''\footnote{\textsc{manguel}, \emph{op.\,cit.}, p.\,34.}. É por meio do fino risco da caneta que a tinta, ao
chocar-se com o papel, transforma-se em lugar.

% \begin{center}\adforn{49}\end{center}